% Generated by `taskgen course examples` from arithmetic.montgomery_mul (seed 1).
% Do not edit: change examples.yaml and regenerate.
\begin{example}[умножение по модулю методом Монтгомери]\label{ex:redc-mul}
Вычислить произведение целых чисел по модулю: \[853194 \cdot 990419\ (\textrm{mod}\ 7600419)\] с помощью метода Монтгомери. Обосновать выбор параметра $R$.

\par\noindent\textit{Решение.} $R = 2^{23} = 8388608$ "--- наименьшая степень двойки, превосходящая модуль $7600419$; $R$ взаимно прост с нечётным модулем, а деление на $R$ "--- сдвиг. $m' = -m^{-1} \bmod R = 4325749$, $R^2 \bmod m = 6451918$ (здесь $m = 7600419$).
 $x = \mathrm{REDC}(ab) = ab R^{-1} \bmod m$, затем $\mathrm{REDC}(x \cdot (R^2 \bmod m)) = ab \bmod m$ (условие $T < mR$ выполнено, так как оба множителя меньше $m$).\par $\mathrm{REDC}(T)$: $u = (T \bmod R)\, m' \bmod R$, $t = (T + u m)/R$, при $t \ge m$ $t \gets t - m$.
\[\begin{array}{l|r|r|r|r}
 & T & u & t & \mathrm{REDC}(T) \\ \hline
a \cdot b & 845019548286 & 936342 & 949098 & 949098 \\
x \cdot (R^2 \bmod m) & 6123502469964 & 4672956 & 4963866 & 4963866 \\
\end{array}\]
\[853194 \cdot 990419 \equiv 4963866\ (\textrm{mod}\ 7600419).\]
\end{example}
